% Slide 06: Exploration loop
\begin{frame}[c]{Continuous exploration loop}
\centering
\begin{tikzpicture}[node distance=0.50cm,every node/.style={font=\scriptsize}]
  \node[greybox,minimum width=1.85cm,minimum height=0.82cm]
    (query) {text query};
  \node[bluebox,right=of query,minimum width=2.10cm,minimum height=0.82cm]
    (search) {semantic\\search};
  \node[mintbox,right=of search,minimum width=2.10cm,minimum height=0.82cm]
    (results) {related\\media};
  \node[redbox,right=of results,minimum width=2.20cm,minimum height=0.82cm]
    (exhibition) {spatial\\exhibition};
  \node[bluebox,right=of exhibition,minimum width=2.10cm,minimum height=0.82cm]
    (select) {select an\\exhibit};

  \foreach \a/\b in {query/search,search/results,results/exhibition,exhibition/select}
    \draw[flow] (\a) -- (\b);

  \draw[flow]
    (select.south) .. controls +(0,-1.05) and +(2.0,-1.05) ..
    node[below=4pt,tinylabel] {use the exhibit as the next query}
    (search.south);
\end{tikzpicture}

\vspace{0.32cm}
\begin{columns}[T,onlytextwidth]
  \begin{column}{0.46\textwidth}
    \textcolor{unimint!80!black}{\large\textbf{One concept}}

    \vspace{0.10cm}
    \small
    The interaction is independent of whether a result is an image, video, or 3D model.
  \end{column}
  \begin{column}{0.46\textwidth}
    \textcolor{unimint!80!black}{\large\textbf{Two entry points}}

    \vspace{0.10cm}
    \small
    Exploration can start from language or continue directly from an existing exhibit.
  \end{column}
\end{columns}
\end{frame}
